% The gap of the tetrahedron triple C = {0,1,2} at (4,3).  Since
% S_C - I_3 = 3I_3 - g_3 g_3^T has spectrum {0,3,3}, the level set
% x^T(S_C - I_3)x = 1 is the circular cylinder of radius 1/sqrt(3) about the
% null line R g_3.  Orthographic projection of R^3 with azimuth 15 and
% elevation 70 degrees, scaled by 1.35cm -- the projection of
% figures/tetrahedron.tex, unaltered, so the two figures are registered:
%   X = 1.35(-0.2588 x_1 + 0.9659 x_2)
%   Y = 1.35(-0.9077 x_1 - 0.2432 x_2 + 0.3420 x_3)
% The viewer lies along e = (0.3303, 0.0885, 0.9397), the cross product of the
% two rows above; a larger e-component is nearer.  Of the four atoms, g_0 sits
% at e-component +1.53 of its distance from the axis and so passes in front of
% the cylinder, g_1 and g_2 at -0.52 and -1.01 and pass behind it.
%
% Cylinder frame.  Axis n = g_3/sqrt(3) = (-1,-1,1)/sqrt(3), radius
% r = 1/sqrt(3); an orthonormal basis of n^perp is p = (1,-1,0)/sqrt(2) and
% q = n x p = (1,1,2)/sqrt(6).  A cross-section circle projects to the ellipse
%   (X,Y)(a) = centre + A cos a + B sin a,
%   A = r P(p) = (-0.67498, -0.36623),  B = r P(q) = (0.22500, -0.14857),
% and P(n) = (-0.55113, 1.16360).  The wall faces front where its outward
% normal cos(a) p + sin(a) q has positive e-component, that is for a in
% (-10.33, 169.67) degrees; those two angles are the silhouette, at projected
% offset +-(0.70439, 0.33365) from the axis.  That offset has length
% 0.77941cm = 1.35/sqrt(3), unforeshortened, because a silhouette ruling meets
% the view direction at a right angle; it is perpendicular to P(n) to five
% places.  The two cross-sections drawn sit at n-parameter -2.10 and +2.10,
% clear of the cube: every smaller symmetric choice puts a rim within 0.1cm of
% an atom tip, since g_0 projects to axial coordinate -1.08 and g_3 to +1.73
% while a rim is only 0.234cm deep in that direction.
% The six tetrahedron edges of figures/tetrahedron.tex are dropped: at the line
% weight of the cylinder wall they read as part of it.
\begin{tikzpicture}[
  x=1cm, y=1cm,
  vec/.style={-{Latex[length=2mm,width=1.5mm]}, line width=0.75pt},
  vecomit/.style={-{Latex[length=2mm,width=1.5mm,fill=white]},
                  line width=0.75pt},
  cube/.style={densely dotted, line width=0.3pt, black!35},
  wall/.style={line width=0.5pt, black!70},
  wallback/.style={line width=0.5pt, black!45, dashed,
                   dash pattern=on 1.4pt off 1.2pt},
  nullline/.style={line width=0.5pt, black!55, dashed,
                   dash pattern=on 2.2pt off 1.6pt},
  every node/.style={font=\small}]

% ---- the eight vertices of [-1,1]^3 -------------------------------------
\coordinate (ppp) at ( 0.9546,-1.0920);   % g_0 = ( 1, 1, 1)
\coordinate (ppm) at ( 0.9546,-2.0154);   %       ( 1, 1,-1)
\coordinate (pmp) at (-1.6534,-0.4354);   %       ( 1,-1, 1)
\coordinate (pmm) at (-1.6534,-1.3588);   % g_1 = ( 1,-1,-1)
\coordinate (mpp) at ( 1.6534, 1.3588);   %       (-1, 1, 1)
\coordinate (mpm) at ( 1.6534, 0.4354);   % g_2 = (-1, 1,-1)
\coordinate (mmp) at (-0.9546, 2.0154);   % g_3 = (-1,-1, 1)
\coordinate (mmm) at (-0.9546, 1.0920);   %       (-1,-1,-1)
\coordinate (O)   at ( 0,0);

% ---- the cube, as a frame of reference ----------------------------------
\draw[cube] (ppp)--(ppm) (ppp)--(pmp) (ppp)--(mpp)
            (ppm)--(pmm) (ppm)--(mpm) (pmp)--(pmm) (pmp)--(mmp)
            (mpp)--(mpm) (mpp)--(mmp) (pmm)--(mmm) (mpm)--(mmm) (mmp)--(mmm);

% ---- the null line R g_3, which is the axis ------------------------------
% n-parameter -2.60 and +2.60; the cube vertices (ppm) and (mmp) sit at -+1.7321
\draw[nullline] ( 1.43294,-3.02536)
  node[anchor=west, xshift=2pt, font=\footnotesize] {$\R g_3$}
  -- (-1.43294, 3.02536);

% ---- the two cross-sections of the cylinder -----------------------------
% centres at n-parameter -2.10 and +2.10; back arc first, then the front arc
\foreach \cx/\cy in {1.15737/-2.44356, -1.15737/2.44356}{%
  \draw[wallback] plot[domain=169.67:349.67, samples=61, smooth, variable=\ang]
    ({\cx-0.67498*cos(\ang)+0.22500*sin(\ang)},
     {\cy-0.36623*cos(\ang)-0.14857*sin(\ang)});
  \draw[wall] plot[domain=-10.33:169.67, samples=61, smooth, variable=\ang]
    ({\cx-0.67498*cos(\ang)+0.22500*sin(\ang)},
     {\cy-0.36623*cos(\ang)-0.14857*sin(\ang)});}

% ---- the two silhouette rulings -----------------------------------------
% each centre -+ the silhouette offset, at a = -10.33 and a = 169.67
\draw[wall] ( 0.45299,-2.77721) -- (-1.86176, 2.10991);
\draw[wall] ( 1.86176,-2.10991) -- (-0.45299, 2.77721);

% ---- the four vectors of the design; g_3 is the one C omits ---------------
\draw[vec]     (O)--(ppp) node[below right=-1pt] {$g_0$};
\draw[vec]     (O)--(pmm) node[below left=-1pt]  {$g_1$};
\draw[vec]     (O)--(mpm) node[right=1pt]        {$g_2$};
\draw[vecomit] (O)--(mmp)
  node[anchor=east, xshift=-2pt, yshift=-4pt]   {$g_3$};
\fill (O) circle (1.1pt);

\node[anchor=north, font=\footnotesize] at (0,-3.30)
  {$x\tp(\SC-I_3)x=1$};

% ---- the cross-section, seen along the axis ------------------------------
\begin{scope}[shift={(4.05,0)}]
  \node[anchor=south, font=\footnotesize] at (0,0.86) {$(\R g_3)^{\perp}$};
  \draw[wall] (0,0) circle (0.77941);        % radius 1.35/sqrt(3)
  \draw[wall] (0.13,0) -- (0.77941,0)
    node[anchor=west, xshift=2pt, font=\footnotesize] {$1/\sqrt3$};
  % the axis, out of the page
  \draw[line width=0.5pt] (0,0) circle (0.13);
  \fill (0,0) circle (1.0pt);
  \node[anchor=east, xshift=-2pt, font=\footnotesize] at (-0.13,0) {$g_3$};
  \node[anchor=north, align=center, font=\footnotesize] at (0,-0.94)
    {$x\tp(\SC-I_3)x=3\lVert x\rVert^{2}$\\[2pt]
     $\operatorname{spec}(\SC-I_3)=\{0,3,3\}$};
\end{scope}

\end{tikzpicture}
